% Section 7.2, from lectures/L07/README.md: the objectives, the evaluation questions, and the next
% lecture.

\section{Review}
\label{c7:sec:review}

\needspace{6\baselineskip}
\subsection*{What you should be able to do now}
After this chapter, you should be able to:
\begin{itemize}
  \item \textbf{Implement \code{readReg()} and \code{writeReg()}} as byte sequences that match the
    protocol exactly, and explain why the four data bytes go most significant first.
  \item \textbf{Explain what \code{const} does and does not reach} through a reference member, and
    justify the \code{const_cast} at the one line where a \code{const} method starts driving
    hardware.
  \item \textbf{Implement the \code{Uart} methods} over that register core, including the poll, read
    and \code{RX_POP} split, injecting the transport through the constructor.
  \item \textbf{Say what each of the three RX accesses would break} if it were omitted, reordered, or
    folded into the one before it.
\end{itemize}

\needspace{6\baselineskip}
\subsection*{Questions to test yourself}
You should be able to explain:
\begin{enumerate}
  \item A \code{writeReg()} sends the four data bytes least significant first: what does the
    peripheral do with the value, and at which point in the book would you first find out?
  \item \code{readReg()} is \code{const} and calls \code{begin()}, \code{transfer()} and
    \code{end()}, which are not. Explain why that compiles with no cast at all, and why the book
    writes the cast anyway.
  \item The read path is poll \code{STATUS}, read \code{RX_DATA}, write \code{RX_POP}. Take each of
    the three away in turn and say what the caller observes.
\end{enumerate}

\needspace{6\baselineskip}
\subsection*{Looking ahead}
In \chapref{8} the driver gets a wire it can be tested against: a scripted
\code{driver::transport::Stub}, the provided host suite that drives the real \code{Uart} over it with
no hardware in sight, and the blocking helpers built on top of the non-blocking core.
